\cotefile{156-4}
\scanpage{58}\mspage{57}
\diagram{\begin{tikzcd}[ampersand replacement=\&,column sep=large]
\underset{\substack{\text{ss-fig de}\\ F}}{\mathcal F_{\leq F}}
 \times
\underset{\substack{\text{ss fig}\\ \text{de }G}}{\mathcal F_{\leq G}}
 \arrow[r,"\sim"] \&
\underset{\substack{\text{ss. fig}\\ \text{de }F\cup G=H}}{\mathcal F_{\leq F\cup G}}
\end{tikzcd}}

(On récupère $F',G'$ \unsure{:|à} partir de $H'=F'\cup G'$ comme $H\cap F$ et $H\cap G$.) \unsure{Item|Idem} pour \unsure{une} \deleted{\unsure{familles|famille} \unsure{mutuellement} disjointes d'une figure (\unsure{y}} \inserted{des figures deux à deux disjointes \unsure{fais.t une|ds une}} configuration (famille \unsure{pas|par} \unsure{néc.t|néct} finie).

\uline{Démonstration} \unsure{Particulièrt|Particulièrmt} claire, quand on interprète $\mathcal F$ en termes de l'ens $\unsure{\mathcal L|L}\unsure{\subseteq|\subset}\mathcal P(M)$ satisfaisant à a) b) c). L'ens ordonné $M_{\leq F\vee G}$ s'identifie à la réunion disjointe dans l'ens ordonné somme de $M_{\leq F}$ et $M_{\leq G}$, donc ses parties fermées \deleted{\unsure{sont}} \unsure{s'en|sont} \unsure{définies|définissent|déduisent} par les couples de parties fermées de $M_{\leq F}$ et $M_{\leq G}$ \ldots

\uline{Remarque}
% slanted
\margin{à joindre à la remarque de la page 55}
En termes de l'ens \uline{ordonné} $M$, la donnée de $\mathcal L$ équivaut \unsure{à|:} un ens de parties \uline{fermées} de $M$, satisfaisant
\[
\begin{aligned}
\text{a) }&\deleted{\unsure{\mathcal L}}\forall\,\deleted{\unsure{=}}A\in\mathcal L,\quad
 \deleted{\text{\unsure{toute}}}\inserted{\text{et}}\ \text{partie fermée }A'\text{ de }A,\ \text{on a }A'\in\mathcal L\\
&\deleted{\text{b) }\varnothing\in\mathcal L\quad
 (\Longleftrightarrow\mathcal L\ne\varnothing,\ \text{\unsure{conséquence} de a), et}}\\
\text{b) }&\bigcup_{A\in\mathcal L} A=M\quad\deleted{\unsure{\forall x}}\quad
 \left[\,\underset{\substack{\text{à cause}\\\text{de a)}}}{\Longleftrightarrow}\
 \forall x\in M,\ M_{\leq x}\in\mathcal L\,\right]\\
\text{c) }&\varnothing\in\mathcal L\quad
 \bigl(\underset{\deleted{\unsure{\text{a) b)}}}}{\Longleftrightarrow}
 \mathcal L\ne\varnothing,\ \text{et \unsure{automatiqt} satisfaite}\ \deleted{\unsure{\mathcal L}}\\
 &\hspace{7em}\text{(à cause de b)) si }M\ne\varnothing\bigr)
\end{aligned}
\]

\uline{NB} Soit $\mathcal P_f(M)$ l'ens des parties fermées de $M$, $\mathcal P_{f'}(\mathcal F)$ de $\mathcal F$ qui sont stables par \unsure{Sup|sup}. Si $\mathcal F'$ et $M'$ se correspondent \inserted{\unsure{resp|ssys}}, 
\diagram{\fbox{\begin{tikzcd}[ampersand replacement=\&,column sep=large]
\mathcal P_f(\unsure{\mathcal F|M\mathcal F}) \arrow[r,shift left=.6ex,"\varphi"] \& \mathcal P_f(M) \arrow[l,shift left=.6ex,"\psi"]
\end{tikzcd}}}
$\varphi(\mathcal F')=\mathcal F'\cap M$, $\psi(M')=\{F\in\mathcal F\mid\overline F\unsure{\subset|\not\subset}M'\}$ alors $\varphi,\psi$ bij. inverses l'une de l'autre
